\newgeometry{left=5mm, right=5mm, top=7.5mm, bottom=7.5mm}
\begin{landscape}
\begin{center}
    \textbf{КАЛЕНДАРНИЙ ПЛАН ВИКОНАННЯ БАКАЛАВРСЬКОЇ РОБОТИ}\\[1em]
    \begin{singlespace}
    \small
    \renewcommand{\arraystretch}{1.1}
    \begin{tabular}{|c|p{14cm}|p{3cm}|p{2.5cm}|p{2cm}|p{2.25cm}|}
        \hline
        № з/п &
        ПЕРЕЛІК РОБІТ &
        Термін виконання &
        \raggedright Дата ознайомлення наукового керівника &
        \raggedright Підпис наукового керівника &
        Примітки \\
        \hline
        1. & Вибір теми, затвердження її на засіданні кафедри та закріплення наукового керівника.\newline
        Узгодження календарного графіка підготовки бакалаврської роботи.\newline
        Ознайомлення студента з критеріями оцінювання бакалаврської роботи &
        & & & \\
        \hline
        2. & Вивчення джерел літератури, матеріалів архівів, періодичних видань, збір та узагальнення фактів, даних &
        & & & \\
        \hline
        3. & Складання плану бакалаврської роботи та узгодження з науковим керівником &
        & & & \\
        \hline
        4. & Написання розділів роботи &
        & & & \\
        \hline
        5. & Проміжний контроль виконання роботи &
        & & & \\
        \hline
        6. & Написання бакалаврської роботи в цілому, ознайомлення з її першим варіантом наукового керівника &
        & & & \\
        \hline
        & \textbf{Розділ 1}\newline
        (постановка проблеми, теоретичні основи, огляд літературних джерел) &
        & & & \\
        \hline
        & \textbf{Розділ 2}\newline
        (проєктування мови) &
        & & & \\
        \hline
        & \textbf{Розділ 3}\newline
        (реалізація компілятора) &
        & & & \\
        \hline
        & \textbf{Розділ 4}\newline
        (тестування та аналіз результатів) &
        & & & \\
        \hline
        7. & Повне завершення написання бакалаврської роботи, оформлення її згідно з вимогами й подання на відгук науковому керівнику &
        & & & \\
        \hline
        8. & Подання бакалаврської роботи для перевірки письмових робіт студентів НаУКМА на відповідність вимогам академічної доброчесності &
        & & & \\
        \hline
        9. & Публічний захист бакалаврської роботи перед екзаменаційною комісією &
        & & & \\
        \hline
    \end{tabular}\\[1em]
    \begin{flushleft}\leftskip=7.5mm{
        Графік узгоджено ``\underline{\hspace{1cm}}'' \underline{\hspace{3cm}} \thesisyear~р.\\
        Науковий керівник: \supervisorfull\\
        Виконавець бакалаврської роботи: \studentfull
    }\end{flushleft}
    \end{singlespace}
\end{center}
\end{landscape}
\clearpage
\restoregeometry
\newpage
